\documentclass[a4paper,11pt]{article}
\usepackage[T2A]{fontenc}
\usepackage[utf8]{inputenc}
\usepackage[russian]{babel}
\usepackage{amsmath,amssymb,mathtools}
\renewcommand{\familydefault}{\sfdefault}
\usepackage{icomma}
\usepackage{geometry}
\geometry{top=23mm,bottom=22mm,left=23mm,right=23mm}
\usepackage{microtype}
\usepackage{tikz}
\usetikzlibrary{calc,arrows.meta,positioning,shapes.geometric}
\usepackage{tabularx,booktabs,array}
\usepackage{enumitem}
\usepackage{hyperref}
\usepackage{parskip}
\usepackage{fancyhdr}
\usepackage{xcolor}
\definecolor{accent}{HTML}{007E80}
\definecolor{darkaccent}{HTML}{00585D}
\definecolor{textgray}{HTML}{34383A}
\definecolor{softgray}{HTML}{6A7274}
\definecolor{forceblue}{HTML}{267399}
\definecolor{gravred}{HTML}{AA4F50}
\hypersetup{colorlinks=true,linkcolor=darkaccent,urlcolor=darkaccent,pdftitle={Чек-лист. Движение по окружности в вертикальной плоскости},pdfauthor={Лёвин Артём Александрович}}
\setlength{\parindent}{0pt}
\setlength{\parskip}{5.5pt}
\setlist[itemize]{leftmargin=17pt,itemsep=2pt,topsep=3pt}
\setlist[enumerate]{leftmargin=20pt,itemsep=4pt,topsep=4pt}
\pagestyle{fancy}
\fancyhf{}
\fancyhead[L]{\small\color{softgray}Движение по окружности}
\fancyhead[R]{\small\color{softgray}Володя Хачатурян}
\fancyfoot[C]{\small\color{softgray}\thepage}
\renewcommand{\headrulewidth}{0.4pt}
\renewcommand{\headrule}{\hbox to\headwidth{\color{accent}\leaders\hrule height \headrulewidth\hfill}}
\setlength{\headheight}{22pt}
\newcommand{\checkitem}{\ensuremath{\square}\hspace{0.35em}}
\newcommand{\acent}{a_{\text{ц}}}
\newcommand{\res}{\textcolor{darkaccent}}
\newcommand{\mysection}[1]{\section{#1}\vspace{-0.2em}}
\begin{document}
% ---- ТИТУЛЬНЫЙ ЛИСТ ----
\thispagestyle{empty}
\begin{center}
\vspace*{1.5cm}
{\Large\color{darkaccent}ФИЗИКА}\par
\vspace{1.7cm}
{\LARGE Чек-лист}\par
\vspace{1.1cm}
{\Large\color{accent}\bfseries Движение по окружности\\[0.12cm]в вертикальной плоскости}\par
\vspace{1.3cm}
{\normalsize Вес тела, реакция опоры\\и сила натяжения нити}\par
\vfill
{\large Володя Хачатурян}\par
\vspace{0.65cm}
{\large 14 ноября 2026 г.}\par
\vspace{1.15cm}
{\small Лёвин Артём Александрович\\[2mm]эксклюзивно для Володи Хачатуряна}\par
\vspace{1.0cm}
\begin{tikzpicture}[x=1cm,y=1cm]
\draw[accent,line width=0.95pt] (0,0) .. controls (1.1,.55) and (2.0,-.55) .. (3.3,0)
 .. controls (4.4,.55) and (5.3,-.55) .. (6.6,0);
\end{tikzpicture}
\end{center}
\clearpage
% ---- ОСНОВЫ ----
\mysection{Главная идея: силы и ускорение}
При движении по окружности \textbf{ускорение направлено к центру} кривизны траектории. Оно существует даже при неизменном модуле скорости: меняется направление вектора скорости.
\[
\boxed{\acent=\frac{v^2}{R}},\qquad
v\;\text{--- скорость (м/с)},\quad R\;\text{--- радиус (м)}.
\]
\textcolor{darkaccent}{\textbf{Правило проекций.}} Сначала выбирают направление оси $y$ и лишь затем расставляют знаки: сонаправленная с осью проекция положительна, противоположная --- отрицательна.
\[
\boxed{\sum F_y=ma_y}\quad\text{(второй закон Ньютона в проекциях).}
\]
\textcolor{darkaccent}{\textbf{Порядок оформления задачи}}
\begin{enumerate}[label=\arabic*.]
\item Нарисуйте тело \textit{в заданной точке} траектории и отметьте центр окружности.
\item Изобразите только действующие на тело силы: $mg$, $N$ или $T$.
\item Проведите ось $y$ и укажите её положительное направление.
\item Направьте $\acent$ \textit{к центру} и запишите $\sum F_y=ma_y$ со знаками.
\item Подставьте $\acent=v^2/R$, найдите величину и проверьте физический смысл.
\end{enumerate}
\mysection{Сила тяжести, вес и контактная сила}
\begin{itemize}
\item \textbf{Сила тяжести} $mg$ действует \textit{на тело} и направлена вертикально вниз.
\item \textbf{Реакция опоры} $N$ действует \textit{на тело} со стороны поверхности, перпендикулярно поверхности.
\item \textbf{Натяжение} $T$ действует \textit{на тело} вдоль натянутой нити, от тела к точке крепления.
\item \textbf{Вес} $P$ --- сила, с которой тело действует на опору или подвес. При единственном контакте с опорой по модулю $P=N$; при подвесе на нити $P=T$.
\end{itemize}
Силы $N$ и $P$ действуют на \textit{разные тела}: это пара сил действия и противодействия. При \textbf{невесомости} контактное действие тела на опору или подвес равно нулю; сила тяжести $mg$ при этом может сохраняться.
\textcolor{darkaccent}{\textbf{Самопроверка}}
\begin{itemize}[label={},leftmargin=0pt]
\item \checkitem На схеме различаю силу тяжести, реакцию опоры и вес.
\item \checkitem Указываю направление оси и ускорения до записи баланса.
\item \checkitem Помню, что $\acent$ --- ускорение, \textit{отдельной силы} $m\acent$ на схеме нет.
\end{itemize}
\clearpage
% ---- МОСТЫ ----
\mysection{Автомобиль на мосту}
В обеих схемах автомобиль находится в точке, где касательная горизонтальна; трением и вертикальными аэродинамическими силами пренебрегаем. Ось $y$ направлена \textbf{вверх}.
\begin{center}
\begin{tikzpicture}[>=Latex,scale=.93, every node/.style={font=\small}]
% left crest
\begin{scope}[shift={(-3.45,0)}]
\node[font=\bfseries\small,text=darkaccent] at (0,2.45) {Верх выпуклого моста};
\draw[softgray,line width=1.05pt] (-2.5,-.55) .. controls (-1.25,1.0) and (1.25,1.0) .. (2.5,-.55);
\filldraw[fill=textgray,draw=textgray] (-.28,.37) rectangle (.28,.64);
\draw[fill=white,draw=textgray] (-.18,.34) circle (.075);
\draw[fill=white,draw=textgray] (.18,.34) circle (.075);
\fill (0,.68) circle (1.6pt);
\draw[->,line width=1.15pt,forceblue] (0,.73)--(0,1.95) node[left] {$N$};
\draw[->,line width=1.15pt,gravred] (-.36,.68)--(-.36,-.54) node[left] {$mg$};
\draw[->,dashed,thick,darkaccent] (.64,.65)--(.64,-.4) node[right] {$\acent$};
\draw[->,black] (1.85,.75)--(1.85,1.46) node[right] {$+y$};
\end{scope}
% right trough
\begin{scope}[shift={(3.45,0)}]
\node[font=\bfseries\small,text=darkaccent] at (0,2.45) {Низ вогнутого участка};
\draw[softgray,line width=1.05pt] (-2.5,.72) .. controls (-1.25,-.85) and (1.25,-.85) .. (2.5,.72);
\filldraw[fill=textgray,draw=textgray] (-.28,-.38) rectangle (.28,-.11);
\draw[fill=white,draw=textgray] (-.18,-.41) circle (.075);
\draw[fill=white,draw=textgray] (.18,-.41) circle (.075);
\fill (0,-.04) circle (1.6pt);
\draw[->,line width=1.15pt,forceblue] (0,0)--(0,1.3) node[left] {$N$};
\draw[->,line width=1.15pt,gravred] (-.44,-.04)--(-.44,-1.15) node[left] {$mg$};
\draw[->,dashed,thick,darkaccent] (.69,0)--(.69,1.18) node[right] {$\acent$};
\draw[->,black] (1.85,-.6)--(1.85,.13) node[right] {$+y$};
\end{scope}
\end{tikzpicture}
\end{center}
\begin{minipage}[t]{.48\linewidth}
\textcolor{darkaccent}{\textbf{Верх: ускорение вниз}}
\[
N-mg=-m\frac{v^2}{R},
\]
\[
\boxed{N=m\left(g-\frac{v^2}{R}\right)}.
\]
С ростом скорости $N$ уменьшается. Формула применима, пока сохраняется контакт ($N\ge0$).
\end{minipage}\hfill
\begin{minipage}[t]{.48\linewidth}
\textcolor{darkaccent}{\textbf{Низ: ускорение вверх}}
\[
N-mg=+m\frac{v^2}{R},
\]
\[
\boxed{N=m\left(g+\frac{v^2}{R}\right)}.
\]
При положительной скорости $N>mg$: нагрузка со стороны автомобиля на дорогу по модулю возрастает.
\end{minipage}
\vspace{0.65em}
\mysection{Условие потери контакта}
На вершине выпуклого моста предельный случай $N=0$:
\[
mg=m\frac{v^2}{R}\quad\Longrightarrow\quad
\boxed{v_{\text{кр}}=\sqrt{gR}}.
\]
При $v=v_{\text{кр}}$ контактная сила обращается в нуль. При большей скорости движение автомобиля \textit{по этой окружности с сохранением контакта} невозможно: автомобиль отделяется от дороги. Отрицательная реакция опоры физически невозможна.
\textit{Важно:} меньшая мгновенная нагрузка от автомобиля на выпуклой вершине сама по себе не определяет выбор конструкции моста; прочность и безопасность зависят и от других факторов.
\clearpage
% ---- НИТЬ ----
\mysection{Ведёрко на нити: верхняя и нижняя точки}
Нить считаем лёгкой и нерастяжимой, натяжение направлено \textbf{вдоль нити от тела к центру}. Сила тяжести в обеих точках направлена вниз. Радиус траектории равен $R$.
\begin{center}
\begin{tikzpicture}[>=Latex,scale=.86,every node/.style={font=\small}]
\begin{scope}[shift={(-3.55,0)}]
\node[text=darkaccent,font=\bfseries\small] at (0,2.08) {Верхняя точка};
\draw[dashed,softgray] (0,0) circle (1.36);
\fill (0,0) circle (2pt) node[right] {центр};
\draw[thick,black] (0,0)--(0,1.36);
\draw[fill=white,thick] (-.24,1.36) rectangle (.24,1.62);
\draw[->,line width=1.15pt,forceblue] (-.16,1.33)--(-.16,.39) node[left] {$T$};
\draw[->,line width=1.15pt,gravred] (.29,1.33)--(.29,.39) node[right] {$mg$};
\draw[->,dashed,darkaccent,thick] (1.55,1.36)--(1.55,.40) node[right] {$\acent$};
\draw[->] (-1.9,1.3)--(-1.9,.52) node[left] {$+y$};
\end{scope}
\begin{scope}[shift={(3.55,0)}]
\node[text=darkaccent,font=\bfseries\small] at (0,2.08) {Нижняя точка};
\draw[dashed,softgray] (0,0) circle (1.36);
\fill (0,0) circle (2pt) node[right] {центр};
\draw[thick,black] (0,0)--(0,-1.36);
\draw[fill=white,thick] (-.24,-1.62) rectangle (.24,-1.36);
\draw[->,line width=1.15pt,forceblue] (-.23,-1.39)--(-.23,-.43) node[left] {$T$};
\draw[->,line width=1.15pt,gravred] (.3,-1.39)--(.3,-2.13) node[right] {$mg$};
\draw[->,dashed,darkaccent,thick] (1.55,-1.4)--(1.55,-.43) node[right] {$\acent$};
\draw[->] (-1.9,-.95)--(-1.9,-1.76) node[left] {$+y$};
\end{scope}
\end{tikzpicture}
\end{center}
\begin{minipage}[t]{.48\linewidth}
\textcolor{darkaccent}{\textbf{Верх: ось вниз}}
\[
T+mg=+m\frac{v^2}{R},
\]
\[
\boxed{T=m\left(\frac{v^2}{R}-g\right)}.
\]
Для натянутой нити требуется $v\ge\sqrt{gR}$. При равенстве $T=0$ --- предельный случай.
\end{minipage}\hfill
\begin{minipage}[t]{.48\linewidth}
\textcolor{darkaccent}{\textbf{Низ: ось вниз}}
\[
mg-T=-m\frac{v^2}{R},
\]
\[
\boxed{T=m\left(g+\frac{v^2}{R}\right)}.
\]
В нижней точке $T>0$. При большой скорости натяжение может превысить прочность нити.
\end{minipage}
\par\vspace{0.65em}
\textbf{Условие невесомости в верхней точке:} $T=0$ для тела на нити, $N=0$ для автомобиля на вершине моста. В обеих задачах получаем $v=\sqrt{gR}$. При этом сила тяжести \textit{остаётся} равной $mg$.

\textbf{Предел прочности.} Если допустимое натяжение $T_{\max}$ задано, то в нижней точке при $T\le T_{\max}$:
\[
\boxed{v\le \sqrt{R\left(\frac{T_{\max}}{m}-g\right)}}
\qquad (T_{\max}\ge mg).
\]
\clearpage
% ---- АЛГОРИТМ НА ОТДЕЛЬНОЙ СТРАНИЦЕ ----
\section*{Алгоритм: задача на вертикальную окружность}
\addcontentsline{toc}{section}{Алгоритм}
\vspace{1.0cm}
\begin{center}
\begin{tikzpicture}[
 >=Latex,
 node distance=8.7mm,
 every node/.style={align=center,font=\small},
 start/.style={ellipse,draw=darkaccent,line width=.75pt,minimum width=4cm,minimum height=8mm},
 flowstep/.style={rectangle,rounded corners=2pt,draw=accent,line width=.75pt,text width=9.0cm,inner sep=6pt},
 decision/.style={diamond,aspect=2.3,draw=darkaccent,text width=5.5cm,inner sep=2pt},
 terminal/.style={rectangle,rounded corners=2pt,draw=softgray,text width=4.2cm,inner sep=5pt},
 arr/.style={->,line width=.75pt,darkaccent}
]
\node[start] (s) {Начало: известна точка траектории};
\node[flowstep,below=of s] (f) {Нарисуйте тело, силы $mg$, $N$ или $T$ и положение центра окружности};
\node[flowstep,below=of f] (a) {Направьте $\acent$ к центру; задайте положительную ось $y$};
\node[flowstep,below=of a] (b) {Запишите проекции: $\sum F_y=ma_y$; учтите знаки всех проекций};
\node[flowstep,below=of b] (c) {Подставьте $\acent=v^2/R$ и выразите $N$ или $T$};
\node[decision,below=11mm of c] (d) {Получено $N\ge0$ или $T\ge0$?};
\node[terminal,below left=12mm and 9mm of d] (yes) {Да: контакт возможен; при нуле --- предельный случай};
\node[terminal,below right=12mm and 9mm of d] (no) {Нет: контакт невозможен; модель движения надо изменить};
\node[start,below=41mm of d] (z) {Проверка и ответ};
\draw[arr] (s)--(f);
\draw[arr] (f)--(a);
\draw[arr] (a)--(b);
\draw[arr] (b)--(c);
\draw[arr] (c)--(d);
\draw[arr] (d.south west) -| node[pos=.21,above left]{да} (yes.north);
\draw[arr] (d.south east) -| node[pos=.21,above right]{нет} (no.north);
\draw[arr] (yes.south) -- ++(0,-.53) -| (z.west);
\draw[arr] (no.south) -- ++(0,-.53) -| (z.east);
\end{tikzpicture}
\end{center}
\clearpage
% ---- КОНТРОЛЬ ЗНАНИЙ ----
\mysection{Типичные ошибки и способы их избежать}
\begin{tabularx}{\linewidth}{@{}p{.31\linewidth}X@{}}
\toprule
\textbf{Ошибка} & \textbf{Как поступать правильно} \\
\midrule
Знак $m\acent$ выбран «на глаз» & Сначала указать ось, затем сравнить с ней направление ускорения к центру. \\
$\acent$ нарисовано как сила & На силовой схеме изображать реальные силы. Ускорение отмечать отдельно. \\
Вес смешан с силой тяжести & $mg$ действует на тело; $P$ действует со стороны тела на опору или подвес. \\
$N<0$ или $T<0$ принят за ответ & Это признак потери контакта или ослабления нити; прежняя модель движения неприменима. \\
В верхней точке $N=mg+mv^2/R$ & Вверху $\acent$ направлено вниз: $N=mg-mv^2/R$. \\
Нить «тянет вниз» в нижней точке & Натяжение направлено вдоль нити \emph{от тела к креплению}, здесь вверх. \\
Забыто квадратное значение скорости & $\acent=v^2/R$, поэтому увеличение скорости вдвое увеличивает $\acent$ вчетверо. \\
\bottomrule
\end{tabularx}
\vspace{.7em}
\mysection{Сравнение четырёх случаев}
\renewcommand{\arraystretch}{1.22}
\begin{tabularx}{\linewidth}{@{}p{.34\linewidth}p{.22\linewidth}X@{}}
\toprule
\textbf{Ситуация} & \textbf{Куда $\acent$} & \textbf{Искомая сила} \\
\midrule
Верх выпуклого моста & вниз & $N=m(g-v^2/R)$ \\
Низ вогнутой дороги & вверх & $N=m(g+v^2/R)$ \\
Верх окружности, нить & вниз & $T=m(v^2/R-g)$ \\
Низ окружности, нить & вверх & $T=m(g+v^2/R)$ \\
\bottomrule
\end{tabularx}
\vspace{.8em}
\mysection{Проверьте себя перед решением}
\begin{itemize}[label={},leftmargin=0pt]
\item \checkitem Объясняю, почему $\acent$ направлено к центру.
\item \checkitem Уверенно выбираю ось и объясняю каждый знак в балансе.
\item \checkitem Различаю силы $mg$, $N$, $T$ и вес $P$.
\item \checkitem Вывожу четыре формулы из второго закона Ньютона.
\item \checkitem Нахожу скорость потери контакта $\sqrt{gR}$ и проверяю ограничения.
\end{itemize}
\clearpage
% ---- ТРЕНИРОВКА ----
\section*{Тренировочный блок: 10 задач}
\textcolor{softgray}{Путь А. Упражнения от базовых к составным. Во всех задачах, если не оговорено иное, $g=10\,\text{м/с}^2$, трением и сопротивлением воздуха пренебречь.}\par
\textbf{В каждом расчётном задании:} сделайте схему сил, укажите направление оси, запишите баланс и только затем вычисляйте.
\begin{enumerate}[label=\textcolor{darkaccent}{\textbf{\arabic*.}},leftmargin=21pt,itemsep=10pt]
\item Автомобиль проходит вершину выпуклого моста со скоростью $12\,\text{м/с}$. Радиус кривизны $36\,\text{м}$. Найдите центростремительное ускорение и укажите его направление.
\item Автомобиль массой $1000\,\text{кг}$ едет по вершине выпуклого моста: $R=50\,\text{м}$, $v=10\,\text{м/с}$. Найдите силу реакции опоры и вес автомобиля.
\item Автомобиль массой $800\,\text{кг}$ находится в нижней точке вогнутого участка: $R=45\,\text{м}$, $v=15\,\text{м/с}$. Определите реакцию опоры.
\item В верхней точке выпуклого участка автомобиль имеет массу $600\,\text{кг}$, скорость $10\,\text{м/с}$ и радиус кривизны $25\,\text{м}$. Выберите \emph{ось вниз}, запишите баланс и вычислите $N$.
\item Найдите предельную скорость автомобиля на вершине выпуклого моста радиусом $90\,\text{м}$, при которой сила реакции опоры становится равной нулю.
\item Ведёрко массой $2\,\text{кг}$ вращают на нити в вертикальной плоскости. В верхней точке $R=2\,\text{м}$, $v=6\,\text{м/с}$. Найдите натяжение нити.
\item В нижней точке круговой траектории ведёрко массой $3\,\text{кг}$ движется со скоростью $10\,\text{м/с}$; $R=5\,\text{м}$. Найдите натяжение нити, выбрав ось вниз.
\item Автомобиль массой $500\,\text{кг}$ движется через вершину выпуклого моста со скоростью $10\,\text{м/с}$. Реакция опоры равна $3000\,\text{Н}$. Определите радиус кривизны.
\item В нижней точке окружности ведёрко массой $2\,\text{кг}$ вращается на нити длиной $4\,\text{м}$. Максимально допустимое натяжение $92\,\text{Н}$. Найдите наибольшую допустимую скорость в этой точке.
\item Ведёрко массой $1{,}5\,\text{кг}$ движется по окружности радиусом $3\,\text{м}$. Модуль скорости и в верхней, и в нижней точке равен $6\,\text{м/с}$. Вычислите оба натяжения. Выдержит ли нить, если предельное натяжение $30\,\text{Н}$?
\end{enumerate}
\clearpage
% ---- ОТВЕТЫ ----
\section*{Ответы к тренировочному блоку}
\vspace{.8cm}
\renewcommand{\arraystretch}{1.65}
\begin{tabularx}{\linewidth}{@{}p{1.8cm}X@{}}
\toprule
\textbf{№} & \textbf{Ответ} \\
\midrule
1 & $4\,\text{м/с}^2$, вертикально вниз, к центру кривизны. \\
2 & $N=8000\,\text{Н}$; $P=8000\,\text{Н}$. \\
3 & $N=12000\,\text{Н}$. \\
4 & $mg-N=mv^2/R$; $N=3600\,\text{Н}$. \\
5 & $v_{\text{кр}}=30\,\text{м/с}$. \\
6 & $T=16\,\text{Н}$. \\
7 & $mg-T=-mv^2/R$; $T=90\,\text{Н}$. \\
8 & $R=25\,\text{м}$. \\
9 & $v_{\max}=12\,\text{м/с}$. \\
10 & $T_{\text{верх}}=3\,\text{Н}$; $T_{\text{низ}}=33\,\text{Н}$; нить не выдержит нагрузку внизу. \\
\bottomrule
\end{tabularx}
\vfill
\textcolor{darkaccent}{\textbf{Итоговая памятка:}} вся задача строится от чертежа и оси к проекциям второго закона Ньютона. Главные ошибки возникают при выборе знаков и при попытке сохранить контакт, когда контактная сила уже должна стать отрицательной.
\vspace{.7cm}
\begin{center}
{\footnotesize\color{softgray}Лёвин Артём Александрович --- эксклюзивно для Володи Хачатуряна}
\end{center}
\end{document}
